% ============================================================
% Chapter 2: State of the Art
%   Section: Automated Fact-Checking
%     Subsection: Language Models as Verifiers
% Included from chapters/02-state-of-the-art/01-automated-fact-checking/automated-fact-checking.tex
% ============================================================

\subsection{Language Models as Verifiers}

Large language models changed the verdict-prediction stage. A
pre-trained model holds a good deal of factual knowledge in its weights
and can answer some factual questions with no evidence at all
\cite{petroni2019lama}; this \emph{closed-book} behaviour is the
baseline against which the contribution of retrieval is measured in
this work (Section~\ref{sec:rq-metrics}). Retrieval-augmented
generation instead conditions the model on documents retrieved for the
question \cite{lewis2020rag}, which is the shape of an evidence-based
verifier: retrieve, then judge the claim against what was retrieved.
More elaborate verifiers decompose a complex claim into simpler
questions and answer each against evidence \cite{pan2023programfc}.

The same models \emph{hallucinate}: they produce fluent statements that
are unsupported by, or contradict, their input or the world
\cite{ji2023hallucination}. For a verifier this is the central risk:
a verdict whose explanation sounds right and rests on nothing. Three
lines of work address it.

\begin{itemize}
    \item \textbf{Citation.} Asking the model to cite the sources each
    statement rests on makes its grounding checkable; ALCE measures
    citation precision and recall automatically
    \cite{gao2023alce}.
    \item \textbf{Atomic checking.} FActScore splits a generated text
    into atomic facts and checks each against a knowledge source
    \cite{min2023factscore}, the same unit of verification this work
    uses for claims.
    \item \textbf{Consistency and calibration.} SelfCheckGPT detects
    hallucinations from the disagreement of several samples of the same
    model \cite{manakul2023selfcheckgpt}; more generally, the confidence
    a neural model reports is often poorly calibrated
    \cite{guo2017calibration}, so it cannot be taken at face value.
\end{itemize}

The verifier of this work takes the first and the last of these into
its design rather than its evaluation. It reduces the model's task to
a closed classification over five verdicts, requires it to cite the
evidence it used and to quote the passage each source supports or
contradicts, discards any quote that does not appear verbatim in its
source, and recomputes the final confidence from properties of the
evidence instead of trusting the model's own
(Sections~\ref{sec:design-fact-checking} and~\ref{sec:dev-fact-checking}).
